\section{Results and Discussion}

To demonstrate functionality of the software implementation, a 1-second long sample was selected from the `s0306lre' data source.
This waveform contains a single heartbeat pulse, and was shown to yield ideal results with the group LASSO algorithm.
Figure \ref{fig:compression} shows the original selected waveform of the ECG on lead 2, as well as the waveform with compression applied.
Note that a compression ratio of 10:1 is selected, which yields a compressed length of 100 samples.

\begin{figure}[h]
    \centering
    \includegraphics[width=\textwidth]{demo_lead_ii_compression.png}
    \caption{Uncompressed v.s. Compressed ECG Lead II}
    \label{fig:compression}
\end{figure}
\par
Recovering the original signal from the compressed data and sensing matrix yielded the following results, as shown in Figures \ref{fig:lead_ii_recovered} and \ref{fig:all_recovered}. The quality of the recovered signal was measured using the percentage of root-mean-squared difference (PRD) for each lead, which was calculated as shown below \citep{daponte2021}. Table \ref{tab:demo_results} contains the PRD observed for each lead in this example.

\[ PRD = \frac{\left\| x_l - \hat{x}_l \right\|_2}{\left\| x_l \right\|_2} \times 100\% \]

\begin{figure}[H]
    \centering
    \includegraphics[width=.95\textwidth]{demo_lead_ii_recovered.png}
    \caption{Source ECG v.s. Recovered Signal, Lead II}
    \label{fig:lead_ii_recovered}
\end{figure}

\begin{figure}[H]
    \centering
    \includegraphics[width=.95\textwidth]{demo_all_leads_recovered.png}
    \caption{Source ECG v.s. Recovered Signal, All Leads}
    \label{fig:all_recovered}
\end{figure}

\begin{table}[H]
    \centering
    \resizebox{\textwidth}{!}{
        \begin{tabular}{|c|c|c|c|c|c|c|c|c|c|c|c|c|}
            \hline
            Lead & I & II & III & aVR & aVL & aVF & V1 & V2 & V3 & V4 & V5 & V6 \\
            \hline
            PRD & 4.62\% & 0.75\% & 0.78\% & 0.89\% & 0.99\% & 0.74\% & 1.78\% & 1.88\% & 1.26\% & 1.87\% & 1.21\% & 1.46\% \\
            \hline
        \end{tabular}
    }
    \caption{PRD per Lead, CR = 10}
    \label{tab:demo_results}
\end{table}

As demonstrated by the figures, despite applying aggressive compression, each lead of the ECG was able to be recovered with minimal error.
While the PRD was higher on some channels, specifically Lead I, the results were well within the expected range established by the literature \citep{mishra2012, daponte2021}.
Most notably, the recovered P, QRS and T complexes conform exceptionally well to the original signal, maintaining any characteristics necessary for further analysis.
\newline
\par
Extending analysis to the selection of algorithm for recovering the sparse coefficients, a similar methodology was applied to a set of 5 healthy control ECGs from the database.
An arbitrary 1-second long sample was selected from each of the following 5 patient recordings: `s0311lre,' `s0299lre,' `s0308lre,' `s0462\_re' and `s0499\_re.'
The same analysis was then applied at \(CR \in \{5, 10\}\), using both LASSO and Group LASSO to recover the original ECG.
Table \ref{tab:compare_algs} shows the average PRD measured for each algorithm over all trials.

\begin{table}[H]
    \centering
    \begin{tabular}{|c|c|c|c|c|}
        \hline
        \multicolumn{1}{|c|}{} & 
        \multicolumn{2}{c|}{\textbf{CR = 5}} & 
        \multicolumn{2}{c|}{\textbf{CR = 10}} \\
        \hline
        Lead & LASSO & Group LASSO & LASSO & Group LASSO \\
        \hline
        I & 3.86\% & 1.62\% & 9.68\% & 2.37\% \\
        II & 4.85\% & 2.00\% & 10.57\% & 2.90\% \\
        III & 3.14\% & 1.72\% & 9.71\% & 3.50\% \\
        aVR & 3.61\% & 1.42\% & 7.11\% & 1.87\% \\
        aVL & 3.68\% & 1.54\% & 12.49\% & 2.69\% \\
        aVF & 3.96\% & 1.88\% & 12.45\% & 3.79\% \\
        V1 & 4.65\% & 1.78\% & 14.50\% & 2.67\% \\
        V2 & 2.70\% & 1.30\% & 9.10\% & 2.56\% \\
        V3 & 3.49\% & 1.38\% & 10.62\% & 2.46\% \\
        V4 & 2.53\% & 1.29\% & 4.58\% & 2.34\% \\
        V5 & 3.05\% & 1.25\% & 9.63\% & 2.32\% \\
        V6 & 3.11\% & 1.15\% & 8.18\% & 2.04\% \\
        \hline
        All & 3.55\% & 1.52\% & 9.88\% & 2.62\% \\
        \hline
    \end{tabular}
    \caption{Average PRD per Lead, LASSO v.s. Group LASSO}
    \label{tab:compare_algs}
\end{table}

As expected, the Group LASSO implementation performs significantly better, especially at higher compression ratios.
Considering that the LASSO algorithm solely optimizes for general sparsity, it is clear that this algorithm is not the best selection for problems with highly-correlated columns.
Although similar in implementation, Group LASSO proves to be much better suited for this problem as it optimizes for row-sparse solutions.
It should be noted, however, that these results are still lower than what can be achieved by M-FOCUSS as shown by Daponte et al. \citep{daponte2021}.
Although originally considered for this study, M-FOCUSS is not a convex problem and thus proved to be non-trivial to implement.
